\section{Qualitative Predictions}


本稿の枠組みは、意識を単一の量として測定する理論ではない。

したがって、ここで提示する予測は「意識量」の数値予測ではなく、速度差、結合粘性、内部遅延、意味勾配、有効トルク様量が変化したとき、意識状態がどのように変形するかについての定性的予測である。

中心となる関係は、前節で導入した次の有効量である。

\begin{equation}
\mathcal{T}_{\Phi}
\sim
\kappa
\frac{\Delta v}{\mu}
\left|\nabla \Phi\right|
\end{equation}

ここで $\mathcal{T}_{\Phi}$ は、速度スケール差と意味勾配が結合した有効トルク様量である。

この量は、意識を生成するものではない。むしろ、ログ参照がどの程度強く、どの方向へ開かれやすいかを示す構造的指標である。

\subsection{Synchronization Collapse}

この予測は、速度差が小さすぎる場合に関するものである。

\begin{equation}
\Delta v \rightarrow 0
\end{equation}

このとき、高速層と低速層のあいだの境界は同期に潰れる。

処理は安定するが、参照のための時間的厚みが失われる。そのため、意識的アクセスは弱くなり、行為は習慣化・自動化・反射化しやすくなる。

この状態では、処理は起きていても、「いま参照している」という感覚は薄くなる。

したがって、過度な同期は意識を強めるのではなく、むしろ意識的参照を不要にする方向へ働く。

\subsection{Dissociation by Excessive Speed Differential}

この予測は、速度差が大きすぎる場合に関するものである。

\begin{equation}
|\Delta v| \gg 0
\end{equation}

このとき、高速層と低速層の対応関係は維持されにくくなる。

高速層では予測や連想が走るが、それが低速層のログとして安定して沈殿しない。その結果、意識的アクセスは断片化し、夢、白昼夢、解離、幻覚様状態に近い構造が現れやすくなる。

この状態では、処理の強度は高くても、ログ参照の一貫性は低下する。

つまり、意識が強くなるのではなく、参照が散乱する。

\subsection{Coupling Viscosity and State Transition}

この予測は、結合粘性 $\mu$ の変化に関するものである。

$\mu$ が低すぎる場合、速度差は十分に滑らず、層どうしが癒着しやすくなる。

\begin{equation}
\mu \rightarrow 0
\end{equation}

この状態では、処理は硬直し、特定のログや意味井戸に固定されやすい。

一方で、$\mu$ が高すぎる場合、層どうしは滑りすぎる。

\begin{equation}
\mu \rightarrow \infty
\end{equation}

この状態では、高速層の変化が低速層に十分伝達されず、ログ参照は不安定化する。

したがって、意識状態は $\mu$ の中間領域においてもっとも安定する。

\begin{equation}
\mu_{\min} < \mu < \mu_{\max}
\end{equation}

この予測は、意識状態を相図的に捉える可能性を示している。

すなわち、意識は ON/OFF の二値ではなく、結合粘性と速度差によって連続的に遷移するレジームである。

\subsection{Delay Window}

この予測は、内部遅延 $\tau$ に関するものである。

\begin{equation}
\tau \propto \frac{\Delta v}{\mu}
\end{equation}

$\tau$ が小さすぎる場合、処理は即時反応に近づき、低速層へ沈殿する前に消える。

\begin{equation}
\tau < \tau_{\min}
\end{equation}

この状態では、意識的参照は弱く、行動は自動処理に近づく。

逆に、$\tau$ が大きすぎる場合、現在処理とログ沈殿の対応が崩れる。

\begin{equation}
\tau > \tau_{\max}
\end{equation}

この状態では、意識は遅延し、断片化し、自己感や時間感覚の不安定化が起こりやすくなる。

したがって、安定した意識には遅延窓が必要である。

\begin{equation}
\tau_{\min} < \tau < \tau_{\max}
\end{equation}

\subsection{Meaning Gradient Bias}

この予測は、意味勾配 $\nabla \Phi$ に関するものである。

ログ参照は、すべてのログに対して均等に開かれるわけではない。

情報ポテンシャル $\Phi$ の勾配が大きい領域では、処理はその方向へ引き寄せられやすい。

\begin{equation}
|\nabla \Phi| \uparrow
\quad \Rightarrow \quad
\text{referencing bias} \uparrow
\end{equation}

このため、強い意味勾配をもつ記憶、感情、概念、物語は、意識的参照に入りやすくなる。

ただし、意味勾配が強すぎる場合、参照は特定の井戸へ固定される。

この状態は、反復思考、強迫的注意、固着、トラウマ的再参照などとして現れうる。

したがって、意味勾配は意識を方向づけるが、過剰な勾配は参照の自由度を奪う。

\subsection{Qualia and Measurement}

この予測は、クオリアと測定に関するものである。

クオリアが高速層と低速層の境界現象であるなら、測定はその境界を変更する。

測定は、ライブな境界状態を低速で報告可能なログへ変換する。

したがって、測定精度を上げるほど、観測されるのはクオリアそのものではなく、クオリアが沈殿した後の残余になる。

\begin{equation}
\text{measurement}
\quad \Rightarrow \quad
\text{boundary closure}
\quad \Rightarrow \quad
\text{sedimented residue}
\end{equation}

この予測は、クオリアが測れないという単純な主張ではない。

測定可能なのは、クオリアそのものではなく、クオリアが通過した後に残るログ構造である、という主張である。

\subsection{Measurement-Induced Residue Shift}

測定要求が強くなるほど、ライブなクオリアは低速層へ沈殿し、報告可能な残余として安定化する。

そのため、測定精度の上昇は、クオリアそのものの捕捉ではなく、沈殿残余の明瞭化として現れる。

この予測は、測定を単なる観察ではなく、意識状態のレジームを変える操作として扱うことを要求する。

同じ経験に対しても、報告要求、選択肢、言語化、数値評価の強さが変われば、残るログ構造は変化する。

したがって、実験で得られるクオリア報告は、ライブな境界状態そのものではなく、測定条件によって形成された residue として解釈される必要がある。

\subsection{Split-Brain Asymmetry}

この予測は、分離脳実験に関するものである。

分離脳において意識が二重化するのではない。

予測されるのは、ログ参照経路の非対称化である。

処理が一方の経路で起きても、それが低速で報告可能なログへ到達できなければ、明示的意識には入らない。

したがって、分離脳で観測される合理化は、別の自己が語っているのではなく、参照できないログの欠落を低速層が補完している状態として説明される。

\subsection{Non-Local Regime Prediction}

意識状態が単一部位ではなく時間的レジームであるなら、同じ局所活動でも、速度差、遅延、意味勾配、報告条件が変われば、意識報告は変化する。

したがって、神経相関は単独では意識状態を決定せず、時間的結合条件と併せて解釈される必要がある。

この予測は、神経活動の局在的対応を否定するものではない。

むしろ、局所活動は consciousness window の中でどの時間的結合に参加しているかによって、異なる意識的意味を持つという主張である。

したがって、同一の部位活動でも、低速ログへの沈殿、報告経路への接続、意味勾配の強さが異なれば、主観報告と行動反応は異なるはずである。

\subsection{Consciousness Window as Phase Diagram}

以上の予測は、意識状態が一点ではなく窓として存在することを示す。

\begin{equation}
0 < |\Delta v| < \Delta v_{\max}
\end{equation}

\begin{equation}
\mu_{\min} < \mu < \mu_{\max}
\end{equation}

\begin{equation}
\tau_{\min} < \tau < \tau_{\max}
\end{equation}

\begin{equation}
\mathcal{T}_{\min}
<
\mathcal{T}_{\Phi}
<
\mathcal{T}_{\max}
\end{equation}

この consciousness window の内側では、処理は完全同期せず、完全分離もせず、ログ参照が安定して開かれる。

この窓は、相図として読むことができる。

横軸に速度差 $|\Delta v|$、縦軸に結合粘性 $\mu$ を置けば、同期崩壊、分離、固着、断片化、安定ログ参照は、それぞれ異なるレジームとして配置される。

さらに、内部遅延 $\tau$ と意味勾配 $|\nabla \Phi|$ は、この相図上の状態を厚み方向へ拡張する制御変数として働く。

したがって、consciousness window は ON/OFF の境界線ではない。

それは、複数の制御変数によって形成される有限の安定領域である。

窓の外側では、意識状態はそれぞれ異なる方向に崩れる。

\begin{itemize}
\item 同期側に崩れると、自動化・習慣化・反射化が強まる。
\item 分離側に崩れると、断片化・夢様状態・解離・幻覚様状態が強まる。
\item 粘性が低すぎると、意味井戸への固定が強まる。
\item 粘性が高すぎると、参照の連続性が失われる。
\item 意味勾配が弱すぎると、参照は方向を失う。
\item 意味勾配が強すぎると、参照は固着する。
\end{itemize}

\subsection{Testable Directions}

本稿は、意識を直接測定可能な単一量として扱わない。

しかし、周辺量の変化から、意識状態の遷移を検証する方向は開かれている。

検証可能な方向として、少なくとも次のものがある。

\begin{enumerate}
\item 神経活動の時間スケール差と主観報告の安定性の対応。
\item 注意操作による $\tau$ の変化。
\item 意味負荷の高い刺激に対する参照バイアスの増大。
\item 測定・報告要求の強化によるクオリア様報告の変質。
\item 睡眠、夢、麻酔、解離状態における速度差・粘性・遅延窓の変化。
\item 言語報告不能な処理におけるログ参照経路の非対称性。
\item 測定要求の強度による residue の明瞭化とライブなクオリア報告の変質。
\item 同一局所活動に対する時間的結合条件の違いによる意識報告の変化。
\item 速度差と結合粘性を軸にした意識状態の相図的分類。
\end{enumerate}

これらはいずれも、本稿の枠組みを完成させるものではない。

むしろ、意識を生成物として探すのではなく、時間スケール間の非閉包状態として調べるための入口である。

\subsection{Summary}

本稿の定性的予測は、次の一点に集約される。

意識状態は、速度差、結合粘性、内部遅延、意味勾配が中間的な有効領域に入ったときに安定する。

\begin{equation}
\text{stable consciousness}
\quad \Leftrightarrow \quad
(\Delta v,\mu,\tau,\mathcal{T}_{\Phi},\nabla\Phi)
\in
\text{consciousness window}
\end{equation}

この窓は固定された場所ではない。

それは、時間的非閉包が破綻せずに維持される動的領域である。

したがって、本稿では意識を生成された対象としてモデル化しない。

意識は、条件が揃ったときに開かれ、条件が崩れたときに閉じる。
